\documentclass[12pt]{article}
\usepackage{custom}

% Data reading
\newcommand{\data}[1]{\input{../output/constants/#1.txt}\unskip}

% External references for any separately compiled online appendix.
\usepackage{xr}
\externaldocument{"manuscript"}

\begin{document}

\title{Online Appendix for ``Broad Validity of the First-Order Approach in Moral Hazard''}
\author{Eduardo M. Azevedo and Ilan Wolff}
\date{\today}
\maketitle

This online appendix derives the functional forms reported in Tables
\ref{tab:utility_functions} and \ref{tab:dists} and provides additional
information on the numerical benchmarks.

\section{Derivations of Scores for Particular Distributions}

For each distribution, the action score is
\[
    S(y\mid a):=\partial_a\log f(y\mid a)
    =\frac{f_a(y\mid a)}{f(y\mid a)}.
\]
Evaluating this expression at the intended action gives the score used in the
body of the paper, $s(y)=S(y\mid a_0)$.

\subsection{Gaussian}

For Gaussian output with mean $a$ and variance $\sigma^2$,
\[
    f(y\mid a)
    =\frac{1}{\sqrt{2\pi\sigma^2}}
      \exp\left(-\frac{(y-a)^2}{2\sigma^2}\right).
\]
Therefore
\[
    S(y\mid a)
    =\partial_a\log f(y\mid a)
    =\frac{y-a}{\sigma^2}.
\]

\subsection{Lognormal}

Suppose $\log y$ is Gaussian with mean $a$ and variance $\sigma^2$.  For
$y>0$,
\[
    f(y\mid a)
    =\frac{1}{y\sqrt{2\pi\sigma^2}}
      \exp\left(-\frac{(\log y-a)^2}{2\sigma^2}\right).
\]
Differentiating its log gives
\[
    S(y\mid a)=\frac{\log y-a}{\sigma^2}.
\]

\subsection{Poisson}

For $a>0$ and $y\in\{0,1,\ldots\}$,
\[
    f(y\mid a)=\frac{a^y e^{-a}}{y!}.
\]
Thus
\[
    S(y\mid a)
    =\partial_a\bigl[y\log a-a-\log(y!)\bigr]
    =\frac{y}{a}-1
    =\frac{y-a}{a}.
\]

\subsection{Exponential}

For the exponential distribution parametrized by its mean $a>0$,
\[
    f(y\mid a)=\frac{1}{a}\exp\left(-\frac{y}{a}\right),
    \qquad y\geq0.
\]
Hence
\[
    S(y\mid a)
    =\partial_a\left[-\log a-\frac{y}{a}\right]
    =-\frac{1}{a}+\frac{y}{a^2}
    =\frac{y-a}{a^2}.
\]

\subsection{Bernoulli}

For $a\in(0,1)$ and $y\in\{0,1\}$,
\[
    f(y\mid a)=a^y(1-a)^{1-y}.
\]
Therefore
\[
    S(y\mid a)
    =\frac{y}{a}-\frac{1-y}{1-a}
    =\frac{y-a}{a-a^2}.
\]

\subsection{Geometric}

Consider the geometric distribution on $\{1,2,\ldots\}$ parametrized by its
mean $a>1$:
\[
    f(y\mid a)
    =\left(1-\frac{1}{a}\right)^{y-1}\frac{1}{a}.
\]
Its log density is
\[
    \log f(y\mid a)
    =(y-1)\bigl[\log(a-1)-\log a\bigr]-\log a.
\]
Consequently,
\[
    S(y\mid a)
    =(y-1)\left(\frac{1}{a-1}-\frac{1}{a}\right)-\frac{1}{a}
    =\frac{y-a}{a^2-a}.
\]

\subsection{Binomial}

For $a\in(0,1)$, a fixed number of trials $n$, and
$y\in\{0,\ldots,n\}$,
\[
    f(y\mid a)=\binom{n}{y}a^y(1-a)^{n-y}.
\]
It follows that
\[
    S(y\mid a)
    =\frac{y}{a}-\frac{n-y}{1-a}
    =\frac{y-na}{a-a^2}.
\]

\subsection{Gamma}

For a fixed shape parameter $n>0$ and scale parameter $a>0$,
\[
    f(y\mid a)
    =\frac{y^{n-1}e^{-y/a}}{\Gamma(n)a^n},
    \qquad y>0.
\]
Thus
\[
    S(y\mid a)
    =\partial_a\left[-\frac{y}{a}-n\log a\right]
    =\frac{y}{a^2}-\frac{n}{a}
    =\frac{y-na}{a^2}.
\]

\subsection{Student's $t$}

For location parameter $a$, scale parameter $\sigma>0$, and degrees of
freedom $\nu>0$,
\[
    f(y\mid a)
    =\frac{\Gamma\left(\frac{\nu+1}{2}\right)}
           {\Gamma\left(\frac{\nu}{2}\right)\sqrt{\pi\nu}\,\sigma}
      \left(1+\frac{(y-a)^2}{\nu\sigma^2}\right)^{-(\nu+1)/2}.
\]
Differentiating the log density gives
\[
    S(y\mid a)
    =\frac{(\nu+1)(y-a)}{\nu\sigma^2+(y-a)^2}.
\]

\subsection{Exponential Family}

Consider
\[
    f(y\mid a)
    =h(y)\exp\bigl(\eta(a)T(y)-A(a)\bigr),
\]
where $h$ is the base density or mass function, $\eta$ is the natural
parameter, $T$ is the sufficient statistic, and $A$ is the log-normalizing
term expressed as a function of action.  Then
\[
    S(y\mid a)
    =\partial_a\log f(y\mid a)
    =\eta'(a)T(y)-A'(a).
\]

\subsection{Additive Model}

Let $Y=a+X$, where $X$ has differentiable density $h$.  Then
\[
    f(y\mid a)=h(y-a),
    \qquad
    f_a(y\mid a)=-h'(y-a).
\]
Therefore
\[
    S(y\mid a)=-\frac{h'(y-a)}{h(y-a)}.
\]

\subsection{Multiplicative Model}

Let $Y=aX$, where $a>0$ and $X$ has differentiable density $h$.  A change of
variables gives
\[
    f(y\mid a)=\frac{1}{a}h\left(\frac{y}{a}\right).
\]
Differentiating with respect to $a$ yields
\[
    f_a(y\mid a)
    =-\frac{1}{a^2}h\left(\frac{y}{a}\right)
     -\frac{y}{a^3}h'\left(\frac{y}{a}\right).
\]
Hence
\[
    S(y\mid a)
    =-\frac{1}{a}
     -\frac{y}{a^2}\frac{h'(y/a)}{h(y/a)}.
\]

\section{Derivations of Links for Particular Utility Functions}

The body of the paper defines utility over total wealth $W$.  Let
\[
    m_0:=\frac{1}{u'(w_0)}.
\]
On the region where limited liability does not bind, the first-order condition
is $u'(W)=1/z$.  The base wealth link and limited-liability wealth link are
therefore
\[
    L_0(z):=(u')^{-1}(1/z),
    \qquad
    L(z):=L_0(z\vee m_0).
\]
For $z(y)=\lambda+\mu s(y)$, the canonical total wealth and wage are
\[
    W_{\lambda,\mu}(y)=L(z(y)),
    \qquad
    w_{\lambda,\mu}(y)=L(z(y))-w_0.
\]
The proofs appendix also uses the corresponding utility link
$\ell(z):=u(L(z))$.  We report it below to distinguish it from the wealth link
$L$ used in the body.

\subsection{Log}

Let
\[
    u(W)=\log W,
    \qquad W>0,
\]
with $w_0>0$.  Since $u'(W)=1/W$, the equation $u'(W)=1/z$ gives
$L_0(z)=z$.  Moreover,
\[
    m_0=\frac{1}{u'(w_0)}=w_0.
\]
Thus
\[
    L(z)=z\vee w_0,
    \qquad
    \ell(z)=\log(z\vee w_0).
\]
The associated canonical wage is
\[
    w_{\lambda,\mu}(y)
    =\bigl[\lambda+\mu s(y)-w_0\bigr]^+.
\]

\subsection{CRRA}

Let
\[
    u(W)=\frac{W^{1-\gamma}}{1-\gamma},
    \qquad W>0,
\]
where $\gamma>0$ and $\gamma\neq1$.  Since $u'(W)=W^{-\gamma}$, the equation
$u'(W)=1/z$ gives
\[
    L_0(z)=z^{1/\gamma}.
\]
Also,
\[
    m_0=\frac{1}{u'(w_0)}=w_0^\gamma.
\]
Consequently,
\[
    L(z)=\bigl(z\vee w_0^\gamma\bigr)^{1/\gamma},
    \qquad
    \ell(z)
    =\frac{\bigl(z\vee w_0^\gamma\bigr)^{(1-\gamma)/\gamma}}
           {1-\gamma}.
\]
The associated canonical wage is
\[
    w_{\lambda,\mu}(y)
    =\Bigl(\bigl[\lambda+\mu s(y)\bigr]\vee w_0^\gamma\Bigr)^{1/\gamma}-w_0.
\]

\subsection{CARA}

Let
\[
    u(W)=-\frac{\exp(-\alpha W)}{\alpha},
\]
where $\alpha>0$.  Since $u'(W)=\exp(-\alpha W)$, the equation
$u'(W)=1/z$ gives
\[
    L_0(z)=\frac{\log z}{\alpha}.
\]
Furthermore,
\[
    m_0=\frac{1}{u'(w_0)}=\exp(\alpha w_0).
\]
It follows that
\[
    L(z)=\frac{\log\bigl(z\vee\exp(\alpha w_0)\bigr)}{\alpha},
    \qquad
    \ell(z)=-\frac{1}{\alpha\bigl(z\vee\exp(\alpha w_0)\bigr)}.
\]
The associated canonical wage is
\[
    w_{\lambda,\mu}(y)
    =\frac{\log\Bigl(\bigl[\lambda+\mu s(y)\bigr]
                       \vee\exp(\alpha w_0)\Bigr)}{\alpha}
     -w_0.
\]

\section{First-Order Approach with CRRA and
\texorpdfstring{$\gamma\in(1/2,1)$}{gamma in (1/2, 1)}}
\label{sec:asymptotic-affinity}

We extend the main theorem to CRRA utility with $\gamma\in(1/2,1)$ under additional assumptions. We use the payment- and utility-contract notation introduced in Appendix~\ref{subsec:appendix-notation}. The main additional restriction is that the expected score is concave in the agent's action. The argument splits the curvature integral into a left tail, which can be ignored asymptotically, and a main region. On the main region, the utility contract is asymptotically affine; by contrast, the main theorem establishes that curvature on this region is small.

\begin{assumption}[CRRA utility]
\label{ass:crra-affinity}
Utility is
\[
    u(W)=\frac{W^{1-\gamma}}{1-\gamma},
    \qquad w_0>0,
    \qquad \gamma\in(1/2,1).
\]
\end{assumption}

Put $p:=(1-\gamma)/\gamma\in(0,1)$.  On the nonbinding region,
\begin{equation}
    \ell_0(z)=\frac{z^p}{1-\gamma},
    \qquad
    \ell_0'(z)=\frac1\gamma z^{p-1}.
    \label{eq:crra-link-affinity}
\end{equation}

\begin{assumption}[Affine-score curvature]
\label{ass:affine-score}
For every $a\in\mathcal A$,
\begin{equation}
    \int s(y)f_{aa}(y\mid a)\,d\nu(y)\leq0.
    \label{eq:affine-score}
\end{equation}
Equivalently, $a\mapsto\E_a[s(Y)]$ is concave.
\end{assumption}

\begin{assumption}[Moments for asymptotic affinity]
\label{ass:affinity-moments}
For some $\eta>0$,
\begin{align*}
 0<\sigma_s^2:=\int s(y)^2f(y\mid a_0)\,d\nu(y)&<\infty,\\
 \int |s(y)|^{2+\eta}f(y\mid a_0)\,d\nu(y)&<\infty,\\
 \sup_{a\in\mathcal A}\int
 (1+|s(y)|^{2+\eta})|f_{aa}(y\mid a)|\,d\nu(y)&<\infty.
\end{align*}
\end{assumption}

\begin{lemma}[Local incentive compatibility multiplier asymptotics]
\label{lem:affinity-multiplier}
Under Assumptions \ref{ass:crra-affinity} and
\ref{ass:affinity-moments},
\[
    \widetilde\mu(\lambda)\ell_0'(\lambda)
    \longrightarrow \frac{c'(a_0)}{\sigma_s^2},
    \qquad
    \frac{\widetilde\mu(\lambda)}{\lambda}\longrightarrow0.
\]
\end{lemma}

\begin{proof}
For fixed $\theta>0$, set
$\mu_\lambda(\theta):=\theta/\ell_0'(\lambda)$.  Since
$\lambda\ell_0'(\lambda)\to\infty$,
$\mu_\lambda(\theta)/\lambda\to0$.  Taylor expansion on
$|s|\leq\lambda/(2\mu_\lambda)$ gives
\[
 \ell(\lambda+\mu_\lambda(\theta)s)
 =\ell(\lambda)+\theta s+r_\lambda(s),
 \qquad
 \int |r_\lambda(s)s|f(y\mid a_0)\,d\nu(y)\to0.
\]
To justify the last assertion, put
$T_\lambda:=\lambda/(2\mu_\lambda)$.  Uniformly for $\theta$ in compact
subsets of $(0,\infty)$, $T_\lambda\asymp\lambda^p$.  On
$\{|s|\leq T_\lambda\}$, Taylor's theorem bounds $|r_\lambda(s)s|$, for
a constant $K$ independent of $\lambda$ and uniform over $\theta$ in the
chosen compact set, by
\[
    K\mu_\lambda^2|\ell_0''(\lambda/2)||s|^3
    \leq K\lambda^{-p}|s|^3.
\]
The $(2+\eta)$ moment makes the integral of this bound
$O(\lambda^{-p\min\{1,\eta\}})$.  Indeed, when $\eta<1$, use
$|s|^3\leq T_\lambda^{1-\eta}|s|^{2+\eta}$ on the central region; when
$\eta\geq1$, the third moment is finite.

On $\{|s|>T_\lambda\}$, sublinear growth gives
\[
 |r_\lambda(s)s|
 \leq K\left[(1+\lambda^p)|s|
       +\mu_\lambda^p|s|^{p+1}+s^2\right].
\]
The $(2+\eta)$ moment and $T_\lambda\asymp\lambda^p$ imply that the
integrals of all three terms are $O(\lambda^{-p\eta})$.  These estimates prove
the asserted convergence, uniformly for $\theta$ in compact sets.  After
multiplication by $s$ and integration, the constant term vanishes because the
score has mean zero, so
\[
 \int \ell(\lambda+\mu_\lambda(\theta)s(y))s(y)
 f(y\mid a_0)\,d\nu(y)
 \longrightarrow\theta\sigma_s^2.
\]
Monotonicity of the left-hand side in $\mu$ brackets the unique local
incentive compatibility multiplier
between any two choices of $\theta$ surrounding
$c'(a_0)/\sigma_s^2$.  This proves the first limit.  The
second follows after division by
$\lambda\ell_0'(\lambda)\to\infty$.
\end{proof}

\begin{lemma}[Uniform affine residual]
\label{lem:uniform-affine-residual}
Under Assumptions \ref{ass:crra-affinity} and
\ref{ass:affinity-moments},
\begin{align}
 \sup_{a\in\mathcal A}\Bigg|\int
 &\left[\ell(\lambda+\widetilde\mu(\lambda)s(y))-\ell(\lambda)
 -\widetilde\mu(\lambda)\ell_0'(\lambda)s(y)\right]
 \nonumber\\[-2pt]
 &\hspace{35mm}\times f_{aa}(y\mid a)\,d\nu(y)\Bigg|
 \longrightarrow0.
 \label{eq:uniform-affine-residual}
\end{align}
\end{lemma}

\begin{proof}
Abbreviate $\widetilde\mu=\widetilde\mu(\lambda)$ and set
$T_\lambda:=\lambda/(2\widetilde\mu)$.  Lemma
\ref{lem:affinity-multiplier} gives
$T_\lambda\asymp\lambda^p$ and
$\widetilde\mu\ell_0'(\lambda)=O(1)$.  On $\{|s|\leq T_\lambda\}$,
Taylor's theorem and \eqref{eq:crra-link-affinity} bound the absolute residual,
for a constant $K$ independent of $\lambda$ and $a$, by
\[
    K\widetilde\mu^2|\ell_0''(\lambda/2)|s^2
    \leq K\lambda^{-p}s^2.
\]
Its integral against $|f_{aa}(y\mid a)|$ is therefore $O(\lambda^{-p})$
uniformly in $a$.

On $\{|s|>T_\lambda\}$, the bound
\[
    |\ell(\lambda+\widetilde\mu s)|
    \leq K(1+\lambda^p+\widetilde\mu^p|s|^p)
\]
and $\widetilde\mu\ell_0'(\lambda)=O(1)$ bound the absolute residual by
\[
    K\left[1+\lambda^p+\widetilde\mu^p|s|^p+|s|\right].
\]
Using the uniform $(2+\eta)$ weighted moment in Assumption
\ref{ass:affinity-moments}, its tail integral is bounded by a constant times
\[
 (1+\lambda^p)T_\lambda^{-(2+\eta)}
 +\widetilde\mu^pT_\lambda^{-(2+\eta-p)}
 +T_\lambda^{-(1+\eta)}
 =O\bigl(\lambda^{-p(1+\eta)}\bigr),
\]
uniformly in $a$.  Combining the central and tail bounds proves
\eqref{eq:uniform-affine-residual}.
\end{proof}

\begin{proposition}[Asymptotic-affinity extension]
\label{prop:asymptotic-affinity}
Suppose $\underline c''>0$ and Assumptions \ref{ass:crra-affinity},
\ref{ass:affine-score}, and \ref{ass:affinity-moments} hold.  Then
\[
    \sup_{a\in\mathcal A}
    \widehat U_{aa}(v_\lambda,a)
    \leq-\underline c''+o(1).
\]
Consequently, the agent's objective is concave for all sufficiently large
$\lambda$, and the first-order approach is valid at all sufficiently high
feasible reservation utilities.
\end{proposition}

\begin{proof}
Abbreviate $\widetilde\mu=\widetilde\mu(\lambda)$.  Using
$\int f_{aa}(y\mid a)d\nu(y)=0$ to remove $\ell(\lambda)$, Lemma
\ref{lem:uniform-affine-residual} gives, uniformly in $a$,
\begin{align*}
 \widehat U_{aa}(v_\lambda,a)
 ={}&\widetilde\mu\ell_0'(\lambda)
 \int s(y)f_{aa}(y\mid a)\,d\nu(y)-c''(a)+o(1).
\end{align*}
The leading term is nonpositive by Assumption \ref{ass:affine-score}, and
$c''(a)\geq\underline c''$, proving the asserted bound.  Local incentive
compatibility at the interior action $a_0$ and concavity imply global
incentive compatibility.  The
last conclusion then follows from the divergence of the smallest supporting
multiplier in Proposition \ref{prop:relaxed-optimal-contract}.
\end{proof}

\begin{corollary}[Gaussian CRRA extension]
\label{cor:gaussian-crra-extension}
For a Gaussian location family on a compact action set, the moment conditions
in Assumption \ref{ass:affinity-moments} hold and
\[
    \int s(y)f_{aa}(y\mid a)\,d\nu(y)
    =\frac{d^2}{da^2}\E_a[s(Y)]=0.
\]
Hence, subject to the maintained assumptions, Gaussian output with CRRA
utility is covered by Theorem \ref{thm:main} for $\gamma>1$ whenever its
capacity condition holds (in particular, for sufficiently low initial
wealth), by its log branch for $\gamma=1$, and by Proposition
\ref{prop:asymptotic-affinity} for $\gamma\in(1/2,1)$.  No general claim is
made here for $\gamma\leq1/2$.
\end{corollary}

\begin{proof}
For Gaussian location output, $s(y)=(y-a_0)/\sigma^2$ is affine in $y$ and all
of its moments against $f(y\mid a_0)$ and the Gaussian derivatives
$f_{aa}(y\mid a)$ are finite uniformly on compact action sets.  Moreover,
$\E_a[s(Y)]=(a-a_0)/\sigma^2$ is affine in $a$, proving the displayed
identity.  The stated cases now follow from the corresponding results.
\end{proof}

\section{Benchmarking Numerical Solvers}
\label{sec:benchmarks}

This section reports benchmarks comparing Algorithm~\ref{alg:cost-minimization} (``Algorithm 1'') with the convex optimization solver described in Section~\ref{sec:numerical-methods}. \input{../output/specs_description.tex}

We benchmark three cases: (i) Gaussian (easy), with parameters as in Figure~\ref{fig:gaussian-log-pp} where the first-order approach holds; (ii) Gaussian (hard), with a low reservation utility where the first-order approach fails; and (iii) Student-$t$ distribution from Figure \ref{fig:t-log-pp}, where the score is not monotone and there are many binding incentive compatibility constraints. For each case, we solve both the cost minimization problem with intended action $a_0=100$ (CMP) and the principal's problem (PP). Each benchmark reports the median time over 20 runs, after discarding one warmup run. The convex solver uses a grid of 201 points for the outcomes and 101 points for the actions.

Table~\ref{tab:benchmarks} reports the results. Algorithm~1 is substantially faster than the convex solver for the cost minimization problem, typically by a factor of 17--90$\times$. The advantage is especially pronounced for the principal's problem, where Algorithm~1 is 2--200$\times$ faster.

\begin{table}[h]
    \centering
    \captionsetup{font=footnotesize}
    \caption{Benchmark results: median runtime in milliseconds.}
    \label{tab:benchmarks}
    \input{../output/benchmarks_table.tex}
    \caption*{\textit{Note:} CMP = cost minimization problem. PP = principal's problem. Each benchmark is the median of 20 runs. ``k'' denotes thousands of milliseconds.}
\end{table}

\end{document}
